\chapter{Z-order and overlays}
\label{ch:zorder}

\begin{aside}
Most UI is flat. A modal dialog, a tooltip, a context menu ---
these need to sit \emph{above} everything else. The terminal
has no Z buffer, so we paint in the right order.
\end{aside}

\section{Painter's algorithm}

The classic 2-D approach: paint back-to-front. Whatever paints
last wins the cell. \maya{} renders the element tree depth-
first, and overlays are appended after the main tree.

\section{Overlay registration}

An overlay is registered with the renderer:

\begin{lstlisting}
ctx.show_overlay(modal_dialog());
\end{lstlisting}

The renderer keeps a list of active overlays and paints each
after the main tree, in the order they were added.

A modal dialog typically covers the screen with a dim
backdrop and a centred panel. The backdrop is itself an
element with semi-transparent style; for terminals without
alpha, ``dim'' means using a darker background colour for
visible cells.

\section{Dismissing overlays}

An overlay is dismissed by removing it from the list. Common
triggers:

\begin{itemize}[leftmargin=*]
  \item Escape key (binds to the focused overlay).
  \item Click outside the overlay's content rectangle.
  \item A button inside the overlay (Cancel, OK).
\end{itemize}

The escape-key binding is automatic for modal overlays. The
focus stack pushes the overlay on show and pops on dismiss.

\section{Focus and Z-order}

Overlays steal focus. The element below cannot receive key
events until the overlay closes. \maya{} stacks focus contexts:

\begin{lstlisting}
ctx.push_focus_scope();
ctx.show_overlay(...);
// when overlay closes:
ctx.pop_focus_scope();
\end{lstlisting}

Inside the overlay, tab cycles within the overlay only. The
outer UI is dimmed visually but not interactive.

\section{Tooltips}

A tooltip is a non-modal overlay tied to a target. It does not
take focus; it does not block input. It vanishes on hover-out
or after a timeout.

\begin{lstlisting}
auto btn = button("Save")
         | tooltip("Ctrl-S also saves");
\end{lstlisting}

The tooltip element registers a hover handler; on hover-in it
shows after a delay, on hover-out it hides.

\section{Context menus}

A context menu (right-click) is a modal overlay positioned at
the cursor location. Selecting an item or pressing escape
dismisses it.

\begin{lstlisting}
on_right_click = [&](int x, int y) {
    ctx.show_overlay(
        context_menu(items, x, y)
    );
};
\end{lstlisting}

The menu element renders at the absolute position, clipped to
the screen edges (so a right-click near the bottom flips the
menu upward).

\section{Toast notifications}

A toast is a transient overlay that appears, lingers, fades.
\maya{}'s toast widget animates in from a corner, holds for 3
seconds, animates out. It does not block input.

\section{The overlay stack}

Multiple overlays can stack: a context menu inside a modal
dialog inside a tooltip-bearing main app. \maya{} paints them
in registration order; the latest is on top.

Dismissing the top reveals the next; the focus and key
dispatch follow.

\section{Implementation: portals}

In React-like systems, overlays are ``portals'' --- their
JSX lives in the component tree but their rendering escapes
to a different DOM root. \maya{}'s analogue:

\begin{lstlisting}
auto card = ...
auto popup = portal(
    panel("Hello") | absolute_position(x, y)
);
auto root = column(card, popup);
\end{lstlisting}

The portal node tells the renderer ``my subtree paints last,
not inline''.

\section{Z-stacking for inline images}

Inline images (Chapter~\ref{ch:images}) sit at a fixed Z layer
defined by the terminal. They cannot be overlaid by cell
paint; their placement is at terminal level, above all cells.

If you need a modal dialog \emph{above} an inline image, the
terminal protocol controls that. Kitty's protocol has
Z-index parameters; iTerm2's does not.

\section{Recap}

\begin{recap}
\begin{itemize}[leftmargin=*]
  \item Painter's algorithm: back-to-front, last wins.
  \item Overlays are appended after the main tree.
  \item Modal overlays push focus scope; non-modal don't.
  \item Tooltips, context menus, toasts are special-case
    overlays.
  \item Inline images sit at their own Z layer.
\end{itemize}
\end{recap}

\section*{Workshop}

\begin{enumerate}[leftmargin=*]
  \item Build a confirmation dialog as a modal overlay.
    Bind escape to cancel, enter to confirm.
  \item Add tooltips to your buttons. Tune the show delay.
  \item Stack a context menu inside a modal. Confirm both
    can be dismissed in the right order.
\end{enumerate}
